\documentclass[11pt,a4paper]{article}
\usepackage[utf8]{inputenc}
\usepackage[T1]{fontenc}
\usepackage[spanish,es-noshorthands]{babel}
\usepackage[margin=1.8cm]{geometry}
\usepackage{booktabs,array,enumitem,amsmath}
\pagestyle{empty}
\begin{document}

\begin{center}
{\Large\bfseries 3B. Leer una \emph{scorecard}}\\[2pt]
Econometría II · PEP 1 · práctica en frío · \textbf{15 min} · calculadora, sin pauta
\end{center}

\begin{center}
\renewcommand{\arraystretch}{1.3}
\begin{tabular}{llr|llr}\toprule
Variable & Tramo & Puntos & Variable & Tramo & Puntos\\\midrule
Edad & $<25$ & 40 & Antigüedad & $<1$ año & 43\\
 & 25--39 & 52 & & 1--3 & 53\\
 & 40 o más & 62 & & 3--7 & 65\\
Historial & con atrasos & 20 & & 7 o más & 60\\
 & sin historial & 35 & \multicolumn{3}{l}{Puntos base: 450 · PDO $=20$}\\
 & al día & 50 & \multicolumn{3}{l}{Se aprueba con 600 puntos o más}\\\bottomrule
\end{tabular}
\end{center}

\noindent\textbf{Ana} tiene 30 años, 8 meses en su trabajo y atrasos previos. \textbf{Bruno} tiene 45 años, 4 años en su trabajo y no tiene historial.

\begin{enumerate}[label=\alph*),itemsep=2pt]
\item Puntaje de cada uno. ¿A quién se aprueba?
\item ¿Cuántas veces mayores son los odds bueno:malo de Bruno que los de Ana?
\item ¿Por qué $<25$ años tiene menos puntos? (relaciónalo con el WOE).
\item ¿Qué tiene de raro la antigüedad (65 y luego 60)? ¿Qué harías?
\end{enumerate}

\bigskip
\noindent\textbf{a)}
\begin{center}
\renewcommand{\arraystretch}{1.8}
\begin{tabular}{|l|p{3.2cm}|p{3.2cm}|}\hline
 & \centering Ana & \centering\arraybackslash Bruno\\\hline
Base & & \\\hline
Edad & & \\\hline
Antigüedad & & \\\hline
Historial & & \\\hline
\textbf{Total} & & \\\hline
¿Aprueba? & & \\\hline
\end{tabular}
\end{center}

\noindent\textbf{b)} Diferencia de puntos $=$ \rule{3cm}{0.4pt} \qquad Odds Bruno / odds Ana $=$ \rule{4cm}{0.4pt}\\[8pt]
Fórmula usada: \rule{13.5cm}{0.4pt}\\[8pt]
\rule{15.4cm}{0.4pt}

\bigskip
\noindent\textbf{c)} \rule{15cm}{0.4pt}\\[10pt]
\rule{15.4cm}{0.4pt}\\[10pt]
\rule{15.4cm}{0.4pt}

\bigskip
\noindent\textbf{d)} \rule{15cm}{0.4pt}\\[10pt]
\rule{15.4cm}{0.4pt}\\[10pt]
\rule{15.4cm}{0.4pt}

\vfill
\noindent\small Hora de inicio: \rule{2cm}{0.4pt} \quad Hora de término: \rule{2cm}{0.4pt}
\end{document}
